% Ban thu: chi giu cac cau co trong CD2

%@ id=TT26.caugiay-l3.003 | trang=1 | so=2 | de=II | vitri=CD2.5,C3.9.6 | muc=9 | tt=0 | soat=ok
\begin{bt}
Cho hai biểu thức: $A = \dfrac{\sqrt{x}+1}{x+\sqrt{x}+1}$ và $B = \dfrac{\sqrt{x}+1}{x-1} - \dfrac{x+2}{x\sqrt{x}-1}$ ($x > 0; x \ne 1$).
\begin{cacy}
\item Tính giá trị của $A$ khi $x = 9$.
\item Rút gọn biểu thức $P = B - A$.
\item Tìm giá trị lớn nhất của biểu thức: $Q = \dfrac{2}{P} + \sqrt{x}$.
\end{cacy}
\loigiai{
\begin{cacy}
\item Thay $x = 9$ (tmđk) vào $A$ ta có
\[ A = \dfrac{\sqrt{x}+1}{x+\sqrt{x}+1} = \dfrac{3+1}{9+3+1} = \dfrac{4}{13}. \]
\item Ta có:
\[ P = B - A = \dfrac{\sqrt{x}+1}{x-1} - \dfrac{x+2}{x\sqrt{x}-1} - \dfrac{\sqrt{x}+1}{x+\sqrt{x}+1} \]
\[ P = \dfrac{\sqrt{x}+1}{(\sqrt{x}+1)(\sqrt{x}-1)} - \dfrac{x+2}{\sqrt{x}^3-1} - \dfrac{\sqrt{x}+1}{x+\sqrt{x}+1} \]
\[ P = \dfrac{1}{\sqrt{x}-1} - \dfrac{x+2}{(\sqrt{x}-1)(x+\sqrt{x}+1)} - \dfrac{\sqrt{x}+1}{x+\sqrt{x}+1} \]
\[ P = \dfrac{x+\sqrt{x}+1}{(\sqrt{x}-1)(x+\sqrt{x}+1)} - \dfrac{x+2}{(\sqrt{x}-1)(x+\sqrt{x}+1)} - \dfrac{(\sqrt{x}+1)(\sqrt{x}-1)}{(\sqrt{x}-1)(x+\sqrt{x}+1)} \]
\[ P = \dfrac{x+\sqrt{x}+1-x-2-x+1}{(\sqrt{x}-1)(x+\sqrt{x}+1)} \]
\[ P = \dfrac{\sqrt{x}-x}{(\sqrt{x}-1)(x+\sqrt{x}+1)} \]
\[ P = \dfrac{-\sqrt{x}(\sqrt{x}-1)}{(\sqrt{x}-1)(x+\sqrt{x}+1)} \]
\[ P = \dfrac{-\sqrt{x}}{x+\sqrt{x}+1}. \]
\item Ta có:
\[ Q = \dfrac{2}{P} + \sqrt{x} = 2 : P + \sqrt{x} = 2 : \dfrac{-\sqrt{x}}{x+\sqrt{x}+1} + \sqrt{x} \]
\[ Q = \sqrt{x} - \dfrac{2(x+\sqrt{x}+1)}{\sqrt{x}} \]
\[ Q = \dfrac{x-2(x+\sqrt{x}+1)}{\sqrt{x}} = \dfrac{x-2x-2\sqrt{x}-2}{\sqrt{x}} \]
\[ = \dfrac{-x-2\sqrt{x}-2}{\sqrt{x}} = -\left(\sqrt{x} + \dfrac{2}{\sqrt{x}} + 2\right). \]
Vì $\sqrt{x} > 0 \Rightarrow \dfrac{2}{\sqrt{x}} \ge 0$, chứng minh được bđt: $a+b \ge 2\sqrt{ab}$ nên theo bất đẳng thức trên, ta có:
\[ \sqrt{x} + \dfrac{2}{\sqrt{x}} \ge 2\sqrt{\sqrt{x} \cdot \dfrac{2}{\sqrt{x}}} = 2\sqrt{2} \]
\[ \sqrt{x} + \dfrac{2}{\sqrt{x}} + 2 \ge 2\sqrt{2} + 2 \]
\[ -\left(\sqrt{x} + \dfrac{2}{\sqrt{x}} + 2\right) \le -2\sqrt{2} - 2 \]
\[ Q \le -2\sqrt{2} - 2. \]
Dấu bằng xảy ra khi:
\[ \sqrt{x} = \dfrac{2}{\sqrt{x}} \Leftrightarrow x = 2. \]
Vậy $Q_{\max} = -2\sqrt{2}-2$ khi và chỉ khi $x = 2$.
\end{cacy}
}
\end{bt}
